\documentclass{beamer}

\usepackage[T2A]{fontenc}
\usepackage[utf8]{inputenc}
\usepackage[english,russian]{babel}
\input{settings.tex}


\title{Дискретные системы}
\subtitle{Теория Управления, лекция 6}
\author{Сергей Савин}
\centering
\date{\mydate}



\begin{document}
\maketitle


%\begin{frame}{Content}
%
%\begin{itemize}
%\item Discrete Dynamics
%\item Stability of the Discrete Dynamics
%\item CT-LTI: Analytical solution
%\begin{itemize}
%    \item Matrix exponential
%    \item Analytical solution
%    \item Forced response
%\end{itemize}
%\item Discretization
%\begin{itemize}
%	\item Exact
%	\item Zero-order hold
%\end{itemize}
%\item Read more
%\end{itemize}
%
%\end{frame}




\begin{frame}{Дискретная динамика}
	\begin{flushleft}
		
		Следующая динамическая система называется \emph{дискретной}:
		
		\begin{equation}
			\bo{x}_{i+1} = \bo{A}\bo{x}_i + \bo{B}\bo{u}_i
		\end{equation}
		
		\begin{itemize}
			\item В уравнении нет производных; система описывается конечно-разностным уравнением, а не ОДУ;
			\item Время изменяется дискретными шагами $i = 0,1,2,\dots$;
			\item Такую систему легко моделировать.
		\end{itemize}
		
		\bigskip
		
		Аффинное управление для этой системы можно записать как:
		
		\begin{equation}
			\bo{u}_i = -\bo{K}\bo{x}_i + \bo{u}_i^*
		\end{equation}
		
	\end{flushleft}
\end{frame}

\begin{frame}{Устойчивость дискретной динамики}
	\begin{flushleft}
		
		Рассмотрим устойчивость дискретной динамической системы, где матрица $\bo{A}$ имеет чисто вещественные собственные значения:
		
		\begin{equation}
			\bo{x}_{i+1} = \bo{A}\bo{x}_i
		\end{equation}		
		
		При собственном разложении $\bo{A} = \bo{V} \bo{D} \bo{V}^{-1}$ (где $\bo{D}$ — диагональная матрица с собственными значениями $\lambda_j$ матрицы $\bo{A}$ на диагонали) и введя обозначение $\bo{z}_i = \bo{V}^{-1}\bo{x}_i$, получаем:
		
		\begin{equation}
			\bo{x}_{i+1} = \bo{V} \bo{D} \bo{V}^{-1}\bo{x}_i
		\end{equation}
		\begin{equation}
			\bo{z}_{i+1} =  \bo{D} \bo{z}_i
		\end{equation}
		
		Это означает, что динамика превратилась в систему независимых скалярных уравнений $z_{j, i+1} =  \lambda_j z_{j, i}$. 
		
	\end{flushleft}
\end{frame}

\begin{frame}{Устойчивость дискретной динамики}
	\framesubtitle{Вещественные собственные значения}
	\begin{flushleft}
		
		Рассматривая отдельные уравнения $z_{j, i+1} =  \lambda_j z_{j, i}$, находим, что абсолютное значение скаляров $z_{j}$ сходится к нулю тогда и только тогда, когда $|\lambda_j| < 1$:
		
		\begin{equation}
			\left| \frac{z_{j, i+1}}{z_{j, i}} \right | =   | \lambda_j |
		\end{equation}
		
	\end{flushleft}
\end{frame}


\begin{frame}{Устойчивость дискретной динамики}
	\framesubtitle{Пара комплексных собственных значений}
	\begin{flushleft}
		
		Рассмотрим устойчивость дискретной динамической системы с матрицей $\bo{A}$ размера 2 на 2:
		
		\begin{equation}
			\begin{bmatrix}
				x_{1, i+1} \\ x_{2, i+1}
			\end{bmatrix}
			= 
			\begin{bmatrix}
				\alpha & -\beta \\ \beta & \alpha
			\end{bmatrix}     
			\begin{bmatrix}
				x_{1, i} \\ x_{2, i}
			\end{bmatrix}
		\end{equation}	
		
		Найдем нормы $\begin{bmatrix}
			x_{1, i+1} \\ x_{2, i+1}
		\end{bmatrix}$ и $\begin{bmatrix}
			x_{1, i} \\ x_{2, i}
		\end{bmatrix}$:
		
		\begin{equation}
			\left | \left|
			\begin{bmatrix}
				x_{1, i} \\ x_{2, i}
			\end{bmatrix}
			\right | \right|^2
			= 
			x_{1, i}^2 + x_{2, i}^2
		\end{equation}
		
		\begin{equation}
			\left | \left|
			\begin{bmatrix}
				x_{1, i+1} \\ x_{2, i+1}
			\end{bmatrix}
			\right | \right|^2
			= 
			(\alpha^2 + \beta^2) (x_{1, i}^2 + x_{2, i}^2)
		\end{equation}
		
	\end{flushleft}
\end{frame}

\begin{frame}{Устойчивость дискретной динамики}
	\framesubtitle{Пара комплексных собственных значений}
	\begin{flushleft}
		
		Мы можем найти отношение норм $\begin{bmatrix}
			x_{1, i+1} \\ x_{2, i+1}
		\end{bmatrix}$ и $\begin{bmatrix}
			x_{1, i} \\ x_{2, i}
		\end{bmatrix}$: 
		
		\begin{equation}
			\left | \left|
			\begin{bmatrix}
				x_{1, i+1} \\ x_{2, i+1}
			\end{bmatrix}
			\right | \right|^2 
			/ 
			\left | \left|
			\begin{bmatrix}
				x_{1, i} \\ x_{2, i}
			\end{bmatrix}
			\right | \right|^2
			=
			\frac{(\alpha^2 + \beta^2) (x_{1, i}^2 + x_{2, i}^2)}{x_{1, i}^2 + x_{2, i}^2}
			= 
			\alpha^2 + \beta^2
		\end{equation}
		
		Вспоминая, что собственные значения системы равны $\lambda  = \alpha \pm j\beta$, мы можем переписать выражение выше как:
		
		\begin{equation}
			\left | \left|
			\begin{bmatrix}
				x_{1, i+1} \\ x_{2, i+1}
			\end{bmatrix}
			\right | \right|^2 
			/ 
			\left | \left|
			\begin{bmatrix}
				x_{1, i} \\ x_{2, i}
			\end{bmatrix}
			\right | \right|^2
			= 
			| \lambda |^2
		\end{equation}		
		
		Мы видим, что $\bo{x}$ сходится к нулю тогда и только тогда, когда $|\lambda| < 1$.
		
	\end{flushleft}
\end{frame}

\begin{frame}{Устойчивость дискретной динамики}
	\begin{flushleft}
		
		Общий критерий устойчивости приведен ниже:
		
		\bigskip
		
		\begin{block}{Асимптотическая устойчивость}
			Дискретная система $\bo{x}_{i+1} = \bo{A}\bo{x}_i$ асимптотически устойчива до тех пор, пока собственные значения матрицы $\bo{A}$ меньше 1 по абсолютной величине: $|\lambda_i(\bo{A})| < 1, \; \forall i$.
		\end{block}
		
		\begin{block}{Маргинальная устойчивость}
			Дискретная система $\bo{x}_{i+1} = \bo{A}\bo{x}_i$ маргинально устойчива до тех пор, пока $|\lambda_i(\bo{A})| \leq 1, \; \forall i$ \textcolor{mynormalgray}{и ни одно собственное значение на единичной окружности не связано с нетривиальным жордановым блоком}.
		\end{block}
		
	\end{flushleft}
\end{frame}


\begin{frame}
	
	\centering{\huge Непрерывные LTI-системы: Аналитическое решение}
	
\end{frame}

\begin{frame}{Матричная экспонента}
	\begin{flushleft}
		
		Экспонента $e^a$ определяется как ряд:
		
		\begin{equation}
			e^a =1+a+\frac{1}{2} a^2
			+\frac{1}{6} a^3 + ... = \sum_{n=0}^{\infty} \frac{1}{n!} a^n
		\end{equation}
		
		Матричная экспонента $e^\bo{A}$ определяется как ряд:
		
		\begin{equation}
			e^\bo{A} = \bo{I}+\bo{A}+\frac{1}{2} \bo{A}\bo{A}
			+\frac{1}{6} \bo{A}\bo{A}\bo{A} + ... = \sum_{n=0}^{\infty} \frac{1}{n!} \bo{A}^n
		\end{equation}
		
	\end{flushleft}
\end{frame}

\begin{frame}{Аналитическое решение ОДУ}
	\begin{flushleft}
		
		ОДУ вида $\dot x = a x$ имеет аналитическое решение $x(t) = e^{at} x(0)$.
		
		\bigskip
		
		ОДУ вида $\dot{\bo{x}} = \bo{A}  \bo{x}$ имеет аналитическое решение $\bo{x}(t) = e^{\bo{A}t}  \bo{x}(0)$.
		
		\bigskip
		
		Проверим, что это действительно решение:
		
		\begin{align}
			\bo{x}(t) &= \left( \bo{I}+\bo{A}t+\frac{1}{2} \bo{A}\bo{A}t^2
			+\frac{1}{6} \bo{A}\bo{A}\bo{A}t^3 + ... \right)  \bo{x}(0) &
			\\
			\dot{\bo{x}}(t) &= \left( \bo{A}+\bo{A}\bo{A}t
			+\frac{1}{2}\bo{A}\bo{A}\bo{A}t^2 + ... \right)  \bo{x}(0) 	&		
			\\
			\dot{\bo{x}}(t) &= \bo{A} \left(\bo{I} +\bo{A}t
			+\frac{1}{2}\bo{A}\bo{A}t^2 + ... \right)  \bo{x}(0) 	&		
			\\
			\dot{\bo{x}}(t) &= \bo{A} e^{\bo{A}t}  \bo{x}(0) 	&		 			
			\\
			\dot{\bo{x}}(t) &= \bo{A} \bo{x}(t)  &\ \ \qed
		\end{align}
		
	\end{flushleft}
\end{frame}

\begin{frame}{Вынужденное движение, 1}
	\begin{flushleft}
		
		ОДУ вида $\dot x = a x + bu(t)$ также имеет аналитическое решение. Чтобы его найти, сначала вычислим следующую производную:
		%
		\begin{equation}
			\frac{d}{dt} \left( e^{-at} x(t) \right) = e^{-at} \dot x(t) - a e^{-at} x(t)
		\end{equation}
		
		Умножая $\dot x = a x + bu(t)$ на $e^{-at}$, получаем:
		
		\begin{align}
			e^{-at} \dot x = e^{-at} a x + e^{-at} bu(t)
			\\
			e^{-at} \dot x -  e^{-at} a x = e^{-at} bu(t)
			\\
			\frac{d}{dt} \left( e^{-at} x(t) \right) = e^{-at} bu(t)
			\\
			\int_{0}^{t} \frac{d}{d\tau} \left( e^{-a\tau} x(\tau) \right)  d\tau = 	\int_{0}^{t}  e^{-a\tau} bu(\tau) \ d\tau
		\end{align}
		
	\end{flushleft}
\end{frame}

\begin{frame}{Вынужденное движение, 2}
	\begin{flushleft}
		
		Продолжая вывод:
		
		\begin{align}
			\int_{0}^{t} \frac{d}{d\tau} \left( e^{-a\tau} x(\tau) \right) d\tau = 	\int_{0}^{t}  e^{-a\tau} bu(\tau) \ d\tau
			\\
			e^{-at} x(t) - x(0) = \int_{0}^{t}  e^{-a\tau} bu(\tau) \ d\tau
			\\
			x(t) = e^{at} x(0) + e^{at} \int_{0}^{t}  e^{-a\tau} bu(\tau) \ d\tau
			\\
			x(t) = e^{at} x(0) + \int_{0}^{t}  e^{a(t-\tau)} bu(\tau) \ d\tau
		\end{align}
		
	\end{flushleft}
\end{frame}

\begin{frame}{Вынужденное движение, 3}
	\begin{flushleft}
		
		Уравнение в пространстве состояний $\dot{\bo{x}} = \bo{A}  \bo{x} + \bo{B}  \bo{u}(t)$ также имеет аналитическое решение:
		
		\begin{equation}
			\bo{x}(t) = e^{\bo{A}t}  \bo{x}(0) + 
			\int_{0}^{t} e^{\bo{A}(t-\tau)} \bo{B}  \bo{u}(\tau) \ d\tau
		\end{equation}
		
		Его можно переписать как:
		
		\begin{equation}
			\bo{x}(t) = e^{\bo{A}t}  \bo{x}(0) + 
			e^{\bo{A}t} \int_{0}^{t} e^{-\bo{A}\tau} \bo{B}  \bo{u}(\tau) \ d\tau
		\end{equation}
		
	\end{flushleft}
\end{frame}

\begin{frame}
	
	\centering{\huge Дискретизация}
	
\end{frame}

\begin{frame}{От аналитического решения к дискретной динамике}
	\begin{flushleft}
		
		Имея решение $\bo{x}(t)$, мы можем изучить, как система эволюционирует от $t=0$ до $t=\Delta t$, предполагая, что $\bo{u}(t) = \bo{u}_0 = \text{const}$ на $\ t \in [ 0, \ \Delta t ]$, и обозначая $\bo{x}(0) = \bo{x}_0$ и $\bo{x}(\Delta t) = \bo{x}_1$:
		
		\begin{equation}
			\bo{x}_1 = e^{\bo{A}\Delta t}  \bo{x}_0 + 
			e^{\bo{A} \Delta t} \int_{0}^{\Delta t} e^{-\bo{A}\tau} \bo{B}  \bo{u}_0 \ d \tau
		\end{equation}
		
		Теперь обозначим:
		%
		\begin{align}
			\bar{\bo{A}} = e^{\bo{A}\Delta t}, \ \ \
			\bar{\bo{B}} = e^{\bo{A} \Delta t} \int_{0}^{\Delta t} e^{-\bo{A}\tau} \bo{B} \ d\tau
		\end{align}
		
		Получаем:
		%
		\begin{equation}
			\bo{x}_1 = \bar{\bo{A}}  \bo{x}_0 + \bar{\bo{B}}  \bo{u}_0
		\end{equation}		
		
	\end{flushleft}
\end{frame}

\begin{frame}{Собственные значения непрерывных и дискретных систем}
	\begin{flushleft}
		
		Рассмотрим систему $\dot{x} = \lambda x$ с решением $x(t) = e^{ \lambda t} x(0)$. Дискретизация будет выглядеть так:
		
		\begin{equation}
			x_1 =e^{ \lambda \Delta t}   x_0 
		\end{equation}
		
		Следовательно, собственное значение дискретной системы равно $\lambda_d = e^{ \lambda \Delta t} $. Полагая $\lambda < 0$, получаем:
		
		\begin{itemize}
			\item При $\Delta t \rightarrow 0$ собственное значение дискретной системы $|\lambda_d| \rightarrow 1$;
			
			\item При $\Delta t \gg 0$ собственное значение дискретной системы $|\lambda_d| \rightarrow 0$;
		\end{itemize}
		
	\end{flushleft}
\end{frame}



\begin{frame}{Приближенная дискретизация, 1}
	\begin{flushleft}
		
		Мы можем разложить точное решение $\bar{\bo{A}} = e^{\bo{A}\Delta t}$ через определение матричной экспоненты:
		%
		\begin{equation}
			\bar{\bo{A}} = e^{\bo{A}\Delta t} =  \sum_{n=0}^{\infty} \frac{1}{n!} \bo{A}^n \Delta t^n
		\end{equation}
		
		Отбрасывая члены высшего порядка, получаем:
		%
		\begin{equation}
			\bo{x}_{i + 1} = ( \bo{I} + \bo{A} \Delta t) \bo{x}_i
		\end{equation}
		
		При уменьшении $\Delta t$ члены высшего порядка (квадратичные и выше по $\Delta t$) стремятся к нулю, что повышает точность приближения.
		
		\bigskip
		
		\textcolor{mynormalgray}
		{
			Тот же результат можно получить, заменив производную в $\dot{\bo{x}} = \bo{A}  \bo{x}$ конечно-разностной аппроксимацией:}
		%
		\begin{equation*}
			\textcolor{mynormalgray}
			{\dot {\bo{x}} \approx \frac{1}{\Delta t}(\bo{x}_{i+1} - \bo{x}_i)}
		\end{equation*}
		
	\end{flushleft}
\end{frame}

\begin{frame}{Приближенная дискретизация, 2}
	\begin{flushleft}
		
		Аппроксимируя \emph{дискретную матрицу состояния} $\bar{\bo{A}}$ и \emph{дискретную матрицу управления} $\bar{\bo{B}}$ следующим образом:
		
		\begin{equation}
			\bar{\bo{A}} = \bo{A} \Delta t + \bo{I}
		\end{equation}
		\begin{equation}
			\bar{\bo{B}} = \bo{B} \Delta t
		\end{equation}
		%
		Получаем дискретную динамику:
		
		\begin{equation}
			\bo{x}_{i+1} = \bar{\bo{A}} \bo{x}_i + \bar{\bo{B}} \bo{u}_i
		\end{equation}
		
	\end{flushleft}
\end{frame}




\begin{frame}{Литература}

\begin{itemize}
	
	\item K. Ogata Modern Control Engineering (Chapter 9.4)
	
	\item Dorf \& Bishop, Modern Control Systems (Chapter 3.3 State differential equation)		
	
\item  \bref{http://cse.lab.imtlucca.it/~bemporad/teaching/ac/pdf/04a-TD_sys.pdf}{Automatic Control 1 Discrete-time linear systems}, Prof. Alberto Bemporad, University of Trento

\item \bref{https://web.mit.edu/2.14/www/Handouts/StateSpaceResponse.pdf}{MIT 2.14, State Space Response}

\end{itemize}

\end{frame}



\myqrframe




\begin{frame}{Жордановы блоки и устойчивость, 1}
	\begin{flushleft}
		
		Рассмотрим матрицу $\bo{A} \in \R^{n \times n}$ с одним жордановым блоком размера $n$ (вся матрица представляет собой жорданов блок). Любой жорданов блок можно записать как:
		%
		\begin{equation}
			\bo{A} = \lambda \bo{I} +\bo{N} 
		\end{equation}
		%
		где $\lambda$ — собственное значение, а $\bo{N} $ — нильпотентная матрица, $\bo{N}^n = \bo{0}$.
		
		\bigskip
		
		Решим уравнение $\dot{\bo{x}} = \bo{A}\bo{x}$ с помощью матричной экспоненты:
		%
		\begin{align}
			\bo{x} (t)= e^{(\lambda \bo{I} +\bo{N})t  }= e^{\lambda t} e^{\bo{N}t} \bo{x} (0)=
			\\
			=  e^{\lambda t} \sum_{k=0}^{n-1} \frac{1}{k!} \bo{N}^k t^k \bo{x} (0)
		\end{align}
		
		Если $\lambda = 0$, то решение является полиномиальным:
		%
		\begin{align}
			\bo{x} (t)=  \sum_{k=0}^{n-1} \frac{1}{k!} \bo{N}^k t^k \bo{x} (0)
		\end{align}
		
	\end{flushleft}
\end{frame}

\begin{frame}{Жордановы блоки и устойчивость, 2}
	\begin{flushleft}
		
		\begin{align}
			\bo{x} (t)=  \sum_{k=0}^{n-1} \frac{1}{k!} \bo{N}^k t^k \bo{x} (0)
		\end{align}
		
		Анализируя решение, замечаем:
		
		\begin{itemize}
			\item Если $\lambda < 0$: экспоненциальное затухание доминирует, система устойчива;
			\item Если $\lambda = 0$: полиномиальный рост, система неустойчива;
			\item Если $\lambda > 0$: экспоненциальный рост, неустойчива.
		\end{itemize}
		
	\end{flushleft}
\end{frame}


\end{document}
